%! Author = sriadityavedantam
%! Date = 3/25/26

\lesson{36}{Wednesday 5 November 2025 9:10}{Day 36}

\begin{theorem}{
    If $f,g$ are continuous, then $\max\{f(x),g(x)\}$ is continuous.
}
	\[
        \max(a,b) = \dfrac{a + b + \abs{a - b}}{2}.
    \]
    So
    \[
        \max\{f,g\}(x) = \dfrac{f(x) + g(x) + \abs{f(x) - g(x)}}{2}.
    \]
    Since sums, differences, and absolute values of continuous functions are continuous, $\max\{f,g\}$ is continuous.
    A corollary is that if $f_1,\ldots,f_N$ are continuous, then $\max\{f_1,\ldots,f_N\}$ is continuous.
\end{theorem}

\begin{example}
    If $f_1,f_2,\ldots$ is a sequence of continuous functions with $f_i(x)\leq C$ for all $x$, define
    \[
        (\sup f_i)(x) \coloneqq \sup_{i} f_i(x).
    \]
    This function need not be continuous.
\end{example}

\begin{example}
    Let $f_m : [0,1]\to\rr$ be given by $f_m(x) = x^m$.
    Then
    \[
        \inf_{m\in\n} x^m =
        \begin{cases}
            0, & 0\leq x < 1,\\
            1, & x = 1.
        \end{cases}
    \]
\end{example}

\begin{example}
    What does
    \[
        \lim_{x\to\infty} f(x) = \infty
    \]
    mean?
\end{example}

\begin{problem}{
    Suppose $\abs{r}<1$.
    Prove that
    \[
        \sum_{n=0}^{\infty}(n+1)r^n
    \]
    converges absolutely and equals $\dfrac{1}{(1-r)^2}$.
}
    \begin{align*}
        \lim_{n\to\infty}
        \abs{\dfrac{(n+2)r^{n+1}}{(n+1)r^n}}
        &=
        \lim_{n\to\infty}
        \dfrac{(n+2)\abs{r}}{n+1}
        = \abs{r} < 1.
    \end{align*}
    So the series converges absolutely by the ratio test.
    Also
    \[
        \sum_{n=0}^{\infty}(n+1)r^n
        =
        \sum_{n=0}^{\infty}r^n \cdot \sum_{n=0}^{\infty}r^n
        =
        \left(\dfrac{1}{1-r}\right)^2.
    \]
\end{problem}

\begin{problem}{
    Prove there exists $c\in(1,2)$ such that $c^3 - c^2 = 1$.
}
    Define $f(x) \coloneqq x^3 - x^2 - 1$.
    Then $f$ is continuous.
    \[
        f(1) = -1, \qquad f(2) = 3.
    \]
    By the Intermediate Value Theorem, there exists $c\in(1,2)$ such that $f(c)=0$.
\end{problem}

\begin{problem}{
    Prove that the irrational numbers are dense in $\rr$.
}
    Given $a<b$, define
    \[
        s \coloneqq a + \dfrac{b-a}{\sqrt{2}}.
    \]
    Then $s\in(a,b)$ and $s\notin\q$.
    Hence $(a,b)\setminus\q \neq \varnothing$.
\end{problem}

\begin{problem}{
    Show that
    \[
        \rr \setminus A^{\circ} = \overline{\rr \setminus A}.
    \]
}
    \textbf{($\subseteq$).}
    If $x\notin A^{\circ}$, then for every $\eps>0$,
    \[
        \vee_\eps(x)\not\subset A.
    \]
    So there exists $y\in \vee_\eps(x)$ with $y\notin A$.
    Hence every neighborhood of $x$ intersects $\rr\setminus A$, so $x\in \overline{\rr\setminus A}$.
    \textbf{($\supseteq$).}
    If $x\in \overline{\rr\setminus A}$, then every neighborhood of $x$ intersects $\rr\setminus A$.
    So no neighborhood of $x$ is contained in $A$, hence $x\notin A^{\circ}$.
\end{problem}

\begin{problem}{
    What does
    \[
        \lim_{x\to\infty} f(x) = \infty
    \]
    mean?
}
    It means:
    For every $M > 0$, there exists $R > 0$ such that
    \[
        x > R \implies f(x) > M.
    \]
\end{problem}

\begin{example}
    Let $f$ be continuous on $\rr$ such that
    \[
        \lim_{x\to\infty} f(x) = \infty
        \quad\text{and}\quad
        \lim_{x\to -\infty} f(x) = \infty.
    \]
    Show $f$ attains a global minimum.
\end{example}

\begin{example}
    Suppose $f : [a,b]\to\rr$ is differentiable and $f\uptick$ is continuous.
    Show $f$ is uniformly differentiable.
\end{example}

\begin{example}
    Suppose $f : (a,b)\to\rr$ is differentiable with $f\uptick(c)=0$ and $f^{\prime\prime}(c)>0$.
    Show $c$ is a local minimum.
\end{example}

\begin{example}
    Suppose $f$ is continuous on $[a,b]$ and differentiable on $(a,b)$.
    If $c\in[a,b]$ is a global maximum, show $c=a$, $c=b$, or $f\uptick(c)=0$.
\end{example}